\documentclass[10pt,a4paper]{amsart}
\usepackage[margin=19mm]{geometry}
\usepackage{amsmath,amssymb,mathtools}
\usepackage[scr=rsfs,cal=euler]{mathalfa}
\usepackage{xfrac}
\usepackage[unicode,hidelinks,bookmarks=false]{hyperref}
\newcommand{\ie}{i.e.\ }
\newcommand{\cf}{cf.\ }
\newcommand{\resp}{respectively\ }
\newcommand{\Subsection}{\subsection}
\providecommand{\nobreakdash}{\nobreak\hbox{-}}
\setlength{\emergencystretch}{2em}
\multlinegap0pt
\usepackage[normalem]{ulem}
\usepackage{xcolor}
\usepackage{amssymb}
\newtheorem{theorem}{Theorem}
\title{Total Positivity of Analytic Bases through Symmetric Functions \protect\thanks{This work was partially supported by Spanish research grants PID2022-138569NB-I00 (MCI/AEI) and RED2022-134176-T (MCI/AEI) and by Gobierno de Aragón (E41$\_$23R, S60$\_$23R).}}
\author{Pablo D\'iaz, Esmeralda Mainar}
\date{Source: arXiv:2601.19998v1 (CC BY 4.0)}
\begin{document}
\maketitle

\noindent\emph{Source and licence.} Total Positivity of Analytic Bases through Symmetric Functions \protect\thanks{This work was partially supported by Spanish research grants PID2022-138569NB-I00 (MCI/AEI) and RED2022-134176-T (MCI/AEI) and by Gobierno de Aragón (E41$\_$23R, S60$\_$23R).}. Version arXiv:2601.19998v1 (CC BY 4.0), distributed by arXiv under the Creative Commons Attribution 4.0 International licence, http://creativecommons.org/licenses/by/4.0/, as recorded in the arXiv licence metadata for this exact version and shown on the abstract page https://arxiv.org/abs/2601.19998v1. Full licence notice: \texttt{licenses/arxiv-2601.19998-CC-BY-4.0.txt}.\par\smallskip

\noindent\textit{Editorial scope.} Counted unit: the article's theorem theo1 (source0.tex lines 240--256). The article's complete original proof of it is delivered with it (source0.tex lines 258--303, one \texttt{proof} environment of 46 lines). No mathematics is written, repaired or re-derived: the statement and the proof are the article's own lines, in the article's own notation, and no retained line is substituted or re-flowed.  Retained uncounted as the source-local context the proof needs: lines 180--182.  Nothing was omitted from any retained span, so the package carries no omission record.  No retained line contains a citation, so the wrapper carries no bibliography and no bibliography entry.  No retained line contains a figure environment, an image-inclusion command or drawing code, so no image is delivered and the package declares no asset dependency.  Editorial formatting only, in addition: an AMS \texttt{amsart} preamble that restates the article's own declarations and macros used by the retained lines, namely the environments \texttt{theorem} on separate counters, together with the standard packages those lines need.  The retained environments print in this document as Theorem 1 in order of appearance, whereas the article prints the same items with the numbering of its own full document.  The complete native source of the article as distributed in the arXiv e-print for this exact version is bundled unchanged as \texttt{sources/193497/source0.tex}.
\begin{equation}\label{schurdef}
s_\lambda(x_1, \dots, x_n) := \frac{ p(x_1, \dots, x_n) }{ V(x_1, \ldots, x_n)}.
\end{equation}  

\begin{theorem} \label{theo1}
Let  \( \mathbf{f} = (f_1, \ldots, f_n) \) be a system of analytic functions on \( I \subseteq \mathbb{R} \) and   \( \mathbf{x} = \{ x_1 , \ldots, x_n \} \) in \( I \) satisfying \( x_1 < \cdots < x_n \). Then 
\begin{equation}\label{MWT}
    \det M_{\bf f, \bf x}  {[i-j+1,\dots,i\,|\,1,\dots, j]}
 = V(x_{i-j+1},\dots,x_i)  \sum_{\lambda\in \Lambda}\frac{1}{C_\lambda}W^t_{{\bf f}}{[j,\lambda]}s_{\lambda}(x_{i-j+1},\dots,x_i),
\end{equation}
   and
   \begin{equation}\label{MWTt}
  \det M_{\bf f, \bf x}^t {[i-j+1,\dots,i\,|\,1,\dots, j]}
 = V(x_{1},\dots,x_j) \sum_{\lambda\in \Lambda}\frac{1}{C_\lambda}W^t_{{\bf f}}{[i,\lambda]}s_{\lambda}(x_{1},\dots,x_j),
\end{equation} 
where, $0 < j \leq i \leq n$, and,  for any $\lambda=(\lambda_1,\ldots,\lambda_j) \in \Lambda$,
\begin{equation}\label{eq:CLambda}
    C_\lambda:=\lambda_j!(\lambda_{j-1}+1)!\cdots (\lambda_1+j-1)!.
\end{equation}
     
\end{theorem}

\begin{proof}
The formula \eqref{MWT} is derived using the following equalities, which are justified below. The proof of \eqref{MWTt} follows  in an entirely analogous manner. 
    \begin{eqnarray*}
    &&\det M_{\bf f, \bf x}{[i-j+1,\dots,i\,|\,1,\dots, j]} =\det \left( \sum_{k=0}^\infty \frac{f_m^{(k)}(0)}{k!}x_l^{k}\right)_{i-j+1 \le l \le i; 1\le m \le j} \nonumber \\%    _{l= i-j+1; m=1}^{i,j}\nonumber \\
%        &&\det M_{\bf f}({\bf x}){[i-j+1,\dots,i\,|\,1,\dots, j]}=\det \left( \sum_{k=0}^\infty \frac{f_m^{(k)}(0)}{k!}x_l^k\right)_{i-j+1\leq l \leq i;\,1\leq m\leq j}\nonumber \\
    &\stackrel{^{(1)}}{=}&{  \sum_{r=1}^j \sum_{k_r=0}^\infty \det\left(  \frac{f_m^{(k_m)}(0)}{k_m!}x_l^{k_m}\right)_{l,m}
    =\sum_{r=1}^j \sum_{k_r=0}^\infty   \bigg[\prod_{m=1}^j \frac{f_m^{(k_m)}(0)}{k_m!}\bigg]\det\left( x_l^{k_m}\right)_{l,m} }\\
    &\stackrel{^{(2)}}{=}&   \sum_{k_1>\cdots> k_j=0}^\infty \sum_{\sigma\in S_j} \bigg[\prod_{m=1}^j \frac{f_{m}^{(k_{\sigma(m)})}(0)}{k_{\sigma(m)}!}\bigg]\det\Big(x_l^{k_{\sigma(m)}}\Big)_{l,m}  \nonumber\\
    &\stackrel{^{(3)}}{=}&\sum_{k_1>\cdots> k_j=0}^\infty \sum_{\sigma\in S_j} \bigg[\prod_{m=1}^j \frac{f_m^{(k_{\sigma(m)})}(0)}{   k_{m}!}\bigg]\text{sgn}(\sigma)\det\Big(x_l^{k_m}\Big)_{l,m}\nonumber\\
    &\stackrel{^{(4)}}{=}&\sum_{k_1>\cdots> k_j=0}^\infty  \bigg[\prod_{m=1}^j \frac{1}{k_m!}\bigg]\det\Big(f_r^{(k_s)}(0)\Big)_{1\leq s,r\leq j}\det\Big(x_l^{k_m}\Big)_{l,m}\nonumber \\
    &=&\sum_{k_1>\cdots> k_j=0}^\infty  \bigg[\prod_{m=1}^j \frac{1}{k_m!}\bigg]W_{\bf f}^t{[1,\dots,j|\,k_j+1,\dots,k_1{  +1}]}\det\Big(x_l^{k_m}\Big)_{l,m}\nonumber \\
    &\stackrel{^{(5)}}{=}&\sum_{\substack{ \lambda\in\Lambda\\
    l(\lambda)\leq j}}  \bigg[\prod_{m=1}^j \frac{1}{(\lambda_m+n-m)!}\bigg]W^t_{{\bf f}}{[j,\lambda]}\det\Big(x_l^{\lambda_m+j-m}\Big)_{l,m}\nonumber \\
    &\stackrel{^{(6)}}{=}& V( x_{i-j+1},\dots,x_i) \sum_{\lambda\in \Lambda}\frac{1}{C_\lambda}W^t_{{\bf f}}{[j,\lambda]}s_{\lambda}(x_{i-j+1},\dots,x_i),
\end{eqnarray*}
where, in the indicated identities, we have taken into account the following aspects.
\begin{enumerate}
    \item[(1)] 
Each $k_1,\dots,k_j$ takes the values of the degree of the terms in the McLaurin expansion of $f_1,\ldots, f_j$, respectively.  Then,  basic properties of determinants have been applied. 
    
 %  We have used the letters $k_1,\dots,k_n$ to store the non-negative integers which specify the degree, in each of the $n$ columns, of the terms in the McLaurin expansion. Then, we have applied basic properties of determinants. 
    \item[(2)] If \(k_r = k_s\) for some \(r, s \in \{1, \dots, j\} \), we have  
    \[\det\Big(x_l^{k_m}\Big)_{l,m}=0.\]
 Therefore, the summatory over \(j\)-tuples \((k_1, \dots, k_j)\) reduces to a sum over tuples with distinct entries. Consequently, we can decompose the sum into two parts: the sum over ordered sequences \(k_1 > \cdots > k_j\) and the sum over all permutations of the elements of each sequence expressed as \((k_{\sigma(1)}, k_{\sigma(2)}, \dots, k_{\sigma(j)})\), where \(\sigma \in S_j\).
 
    
    
    \item[(3)] Be aware that, for $\sigma\in S_j$, we have 
    \begin{equation*}
      \prod_{m=1}^j \frac{1}{k_{\sigma(m)}!}= \prod_{m=1}^j \frac{1}{k_m!}\quad \text{ and }\quad \det\Big(x_l^{k_{\sigma(m)}}\Big)_{l,m}=\det\Big(x_l^{k_{m}}\Big)_{l,m}\text{sgn}(\sigma).
    \end{equation*}
    \item[(4)] Note that
    \begin{equation*}
        \sum_{\sigma\in S_j} \prod_{m=1}^j f_m^{(k_{\sigma(m)})}(0)\text{ sgn}(\sigma)=\det\Big(f_r^{(k_s)}(0)\Big)_{1\leq s,r\leq j}.
    \end{equation*}
    \item[(5)] 
After performing the substitution  
\begin{equation*}
    k_m = \lambda_m + j - m, \quad m = 1, \dots, j,
\end{equation*}  
each tuple \(k_1 > \cdots > k_j\) is transformed into a partition \(\lambda=(\lambda_1, \dots, \lambda_j)\) with at most \(j\) parts. It is straightforward to observe that summing over all sequences \(k_1 > \cdots > k_j\) is equivalent to summing over all partitions \(\lambda\in\Lambda\) with \(l(\lambda) \leq j\).  

\item[(6)] We have applied Jacobi's bialternant formula \eqref{schurdef}.  The restriction \(l(\lambda) \leq j\) is unnecessary in the summation since Schur functions of $j$ variables are zero whenever \(l(\lambda) > j\).  
\end{enumerate}
 \qed
\end{proof}

\end{document}
